% Pursuing Stacks (cotes 134-2 to 134-8): archive cote 134-3, site page 13 (PDF page 14), typescript page 13.
% typescript. Draft compared: ps2.tex:569-633.
% AI-transcribed by Claude Opus 5.5 (claude-opus-5-5, adaptive thinking, effort high), Message Batches API, batch ps1 (msgbatch_01Rv5hCaNTtUt4JHx1yqoHPL), 2026-11-12; not yet checked by eye.
\scanpage{134-3/13}
\mspage{13}
Gr-\uline{stacks}, which will have to \xout{tx} be developped at a later stage.

The idea just comes to my mind whether the exactness condition
implying standar\over{t}{d} amalgamated sums, defining the subcategory of stacks
within $\overset{\hand{$\wedge$}}{\underline{B}}$ , cannot be interpreted in terms of some more or less obvious
topology on $\underline{B}$ turning $\underline{B}$ into a site, as the subcategory of correspon-%
ding \uline{sheaves}. This would mean that the category (Gr-stacks) is in fact
a topos, with the host of eategorical information and topological
intuition that goes with such a situation. In this connection, it is
timely to recall that the related categories (Cat) of ``all'' categories,
and (Groupoids) of ``all'' groupoids, are definitely \uline{not} topoi (if my
recollection is correct - it is\hand{'}nt immediately clear to me why they
are not). This seems to suggest that, granting that (Gr-stacks) is
indeed a topos, that this would be a rather special feature of the
structure species of infinite order we are working with (as \deleted{a}\inserted[typed]{one} re\unsure{w}ard so
to say,among a heap of others, for conceptual sophistication !), in
contrast to the categories (Gr-n-stacks)\inserted[typed]{with finite n}, which presumably are not topoi
(For a definition of Gr-n-stacks in terms of Gr-stacks namely
Gr-$\infty$-stacks, see below.\hand{)}

A related question is whether the category (Gr-stacks) is a model
category for the usual homotopy category, the pair of adjoint functors
bonsidered before satisfying \inserted{moreover} the conditions of Quillen's comparison
theorem\hand{(\textsuperscript{13})}.The obvious idea that comes to mind here, in order to define
the model structure on (Gr-stacks), is to take as ``weak equivalences''
the maps which are transformed into weak equivalences by the topological
realization functor (which should be readily expressible in \inserted[typed]{algebraic} terms \xout{as} ..\hand{)},
\over{\unsure{as}}{as} cofibrations the monomorphisms, and defining fibrations by the
Serre-Quillen lifting property with respect to cofibrations which are
weak equivalences (with the expectation that we even get a ``closed mod\hand{el}
category'' in the sense of Quillen). Here it does'nt look too unreasonab
\hand{-}le to expect the same constructions to work in each of the categories
(Gr-n-stacks), n \over{$>$}{$\geqslant$}1, as well as in the categories (n-stacks) without
Gr , which are still to be defined though.

Coming back upon the question of a suitable topology on $\underline{B}$, the
idea that comes to mind immediately of course is to define\deleted{s} covering
families of an object K of $\underline{B}$ i.e. of $\underline{B}_o$ , in terms of the components
which occur in the description of K as an \inserted[typed]{iterated} amalgamated sum of\deleted{\unsure{s}} cells $D_n$.
A quick glance (too quick a glance ?) seems to show this is indeed a
topology, and that the sheaves for this topology are what we expect.
\deleted{This would imply that the \deleted{inclusion} \inserted{composite} functor}
